\chapter{Open Abdominal Aortic Aneurysm Repair}
\section*{Overview}
Open abdominal aortic aneurysm repair replaces or reinforces a weakened section of the abdominal aorta with a synthetic graft.
\section*{Common Instruments and Equipment}
The exact setup varies by technique and setting. Common items may include:
\begin{itemize}[nosep,leftmargin=*]
  \item vascular surgery instrument set and vascular clamps
  \item prosthetic aortic graft
  \item cell-salvage and blood-management equipment as needed
\end{itemize}

\section*{Why It Is Done}
It is performed for selected large, symptomatic, rapidly enlarging, or ruptured abdominal aortic aneurysms when open repair is appropriate.
\section*{How It Is Performed}
Through an abdominal incision, the aorta is clamped, the aneurysmal segment is opened or bypassed, and a graft is sewn between healthy sections of artery.
\section*{Contraindications}
No single contraindication applies to every patient; the treating team assesses whether the procedure is suitable for the individual and the reason for treatment.

\section*{Precautions}
The care team reviews relevant health conditions, medicines, allergies, and procedure-specific risks, and explains any changes needed for the individual.

\section*{Risks and Recovery}
Bleeding, heart attack, kidney injury, bowel or spinal ischemia, infection, graft complications, and death are possible. Recovery is longer than after endovascular repair.
\section*{References}
\begin{enumerate}\item MedlinePlus. Abdominal Aortic Aneurysm Repair. U.S. National Library of Medicine; 2025.\item Current Medical Diagnosis and Treatment. McGraw-Hill; 2025.\end{enumerate}
\clearpage
